\section{Цепи Маркова}

\begin{definition}[Цепь Маркова]
Рассмотрим последовательность случайных величин $S^{(0)}, S^{(1)}, \dots, S^{(n)}$, принимающих значения из конечного множества состояний $\{s_1, s_2, \dots, s_N\}$.

Начальное распределение: $\vec{p}^{(0)} = (p_0^{(0)}, \dots, p_N^{(0)})$.

Матрица переходных вероятностей $P_{N \times N} = \|p_{ij}\|$ удовлетворяет условиям:
\begin{enumerate}
    \item $p_{ij} \ge 0$
    \item $\sum_{j=1}^{N} p_{ij} = 1$ для всех $i$
\end{enumerate}
\end{definition}

\begin{definition}[Марковское свойство]
Вероятность траектории:
$$ P(S^{(0)} = s_{i_0}, S^{(1)} = s_{i_1}, \dots, S^{(n)} = s_{i_n}) = p_{i_0}^{(0)} \cdot p_{i_0 i_1} \cdot p_{i_1 i_2} \cdots p_{i_{n-1} i_n} $$
для строки $s_{i_0}, s_{i_1}, \dots, s_{i_n}$ длины $n+1$.
\end{definition}

\begin{theorem}[Марковское свойство (условная независимость)]
Для цепи Маркова условные вероятности событий при условии промежуточного состояния распадаются на произведение:
$$ P(AC|B) = P(A|B) \cdot P(C|B) $$
где $A = \{S^{(t)} = s_i\}$, $B = \{S^{(r)} = s_j\}$, $C = \{S^{(v)} = s_k\}$ при $t < r < v$.
\end{theorem}

\begin{proof}
Рассмотрим событие $B = \{S^{(r)} = s_j\}$.
$$ P(B) = \sum_{i_0, i_1, \dots, i_{r-1}, i_{r+1}, \dots, i_n} P(S^{(0)} = s_{i_0}, \dots, S^{(r)} = s_j, \dots, S^{(n)} = s_{i_n}) $$
$$ = \sum p_{i_0}^{(0)} p_{i_0 i_1} \cdots p_{i_{r-1} j} \cdot p_{j i_{r+1}} \cdots p_{i_{n-1} i_n} $$

Для события $C = \{S^{(t)} = s_i, S^{(v)} = s_j\}$:
$$ P(C) = \sum p_{i_0}^{(0)} p_{i_0 i_1} \cdots p_{i_{t-1} i} \cdot p_{i i_{t+1}} \cdots p_{i_{v-1} j} \cdot p_{j i_{v+1}} \cdots p_{i_{n-1} i_n} $$

Разделим сумму на две части (до момента $r$ и после):
$$ P(C) = \underbrace{\sum p_{i_0}^{(0)} p_{i_0 i_1} \cdots p_{i_{r-1} i}}_{P(B)} \cdot \underbrace{\sum p_{i i_{r+1}} \cdots p_{i_{n-1} j}}_{P(C|B)} $$

Тогда:
$$ P(AC|B) = \frac{P(ABC)}{P(B)} = \frac{P(C|AB)P(AB)}{P(B)} = P(A|B) \cdot P(C|AB) = P(A|B) \cdot P(C|B) $$
\end{proof}

\begin{statement}
Сумма вероятностей всех траекторий равна 1:
$$ \sum_{i_0, \dots, i_n} P(S^{(k)} = s_{i_k} \text{ для } k \in [1, n]) = \sum_{i_0} p_{i_0}^{(0)} \sum_{i_1, \dots, i_n} p_{i_0 i_1} \cdots p_{i_{n-1} i_n} = 1 $$
\end{statement}